\vspace*{0pt plus 1fill}
\begin{tucsimplesection}[2\bigskipamount]{Abstract}

\bigskip
The integration of automated AI agents into Decentralized Finance (DeFi) introduces potential
for abstracting complex financial operations behind simple natural language commands. To
enable such automation safely, the ERC-4337 Account Abstraction standard has emerged as a
foundational execution layer, providing programmable smart contract wallets with basic
constraints such as session-scoped spending limits and whitelisted contract enforcement.
However, it cannot evaluate the semantic correctness of the contracts that an AI agent targets.
Current AI-Web3 architectures rely on probabilistic Retrieval-Augmented Generation (RAG)
for on-chain data retrieval, leaving automated agents vulnerable to hallucination-induced
misrouting and adversarial prompt injection attacks, possibly resulting in irreversible loss of
user funds. \\ \\
This thesis aims to tackle this issue through a verification architecture that introduces a
deterministic, graph-based pre-execution validation layer between the AI reasoning stage and
the blockchain execution layer. Instead of directly identifying contract addresses, the AI should
only extract the declared transaction intent, while an automatically populated, domain-specific
Blockchain graph would then be used to verify the involved protocols, pools, and tokens
before allowing the transaction. This graph-based resolution layer could then be combined with
an ERC-4337 smart contract wallet at the execution stage for a layered verification architecture.\\ \\
The objective of this thesis is the creation of a solution or the combination of existing
approaches to solve the above-described problem of verifying DeFi transactions in AI agents.
This includes an analysis of the state of the art of existing verification mechanisms, Blockchain
graphs, graph retrieval methods, and other relevant literature. Based on this analysis, a possible
solution should be conceptualized, demonstrated through a prototypical implementation, and a
suitable evaluation conducted based on its effectiveness and its compliance with requirements
extracted through the literature research.
\end{tucsimplesection}
\vspace*{0pt plus 2.5fill}
